% Generated draft: requires mathematical and source review.
\documentclass[11pt]{article}
\usepackage[margin=25mm]{geometry}
\usepackage{fontspec}
\setmainfont{Latin Modern Roman}
\usepackage{amsmath,amssymb,mathtools,unicode-math}
\setmathfont{Latin Modern Math}
\usepackage{enumitem,needspace,xurl,hyperref}
\setlength{\parindent}{0pt}
\setlength{\parskip}{5pt}
\setlength{\emergencystretch}{3em}
\newcommand{\N}{\mathbb{N}}
\newcommand{\Z}{\mathbb{Z}}
\newcommand{\Q}{\mathbb{Q}}
\newcommand{\R}{\mathbb{R}}
\newcommand{\C}{\mathbb{C}}
\newcommand{\F}{\mathbb{F}}
\newcommand{\A}{\mathbb{A}}
\newcommand{\PP}{\mathbb{P}}
\newcommand{\E}{\mathbb{E}}
\newcommand{\eps}{\varepsilon}
\newcommand{\dd}{\,\mathrm{d}}
\begin{document}
\section*{UnsolvedMath ID: 800459}
\textbf{Source problem code:} catalogue entry
\subsection{Definitions and mathematical statement}
Let \((M_i^n,g_i,p_i)\) be pointed Riemannian \(n\)-manifolds satisfying
\[
\operatorname{Ric}_{g_i}\geq -(n-1)g_i,
\qquad
\operatorname{Vol}_{g_i}\bigl(B_1^{g_i}(p_i)\bigr)>v>0
\quad\text{for every }i,
\]
where \(v\) is independent of \(i\). The supplied record assumes convergence \((M_i^n,g_i,p_i)\to(X,d,p)\), but does not specify its meaning. For the formulation below, interpret it as pointed Gromov--Hausdorff convergence; this interpretation requires source verification.

For \(x\in X\), let \(\mathcal T_xX\) denote the collection, up to pointed isometry, of pointed Gromov--Hausdorff limits
\[
(X,r_j^{-1}d,x)\longrightarrow(T,d_T,o),
\qquad r_j>0,\quad r_j\to0.
\]
Such a limit is called a tangent cone at \(x\).

The supplied conjecture asserts that, outside an exceptional set of ``codimension five,'' all tangent cones at each fixed point are homeomorphic. Under the additional interpretation that this means Hausdorff codimension at least five, a precise candidate statement is: there exists \(E\subseteq X\) such that
\[
\dim_{\mathrm H}E\leq n-5
\]
and
\[
\forall x\in X\setminus E\quad
\forall (T_1,d_1,o_1),(T_2,d_2,o_2)\in\mathcal T_xX,
\qquad T_1\cong_{\mathrm{top}}T_2.
\]
Here \(\cong_{\mathrm{top}}\) means that the underlying topological spaces are homeomorphic. The comparison is between tangent cones at the same point; no comparison between different points is asserted. With the convention \(\dim_{\mathrm H}\varnothing=-\infty\), this candidate interpretation requires \(E=\varnothing\) when \(n<5\).

The Hausdorff-dimension interpretation, the low-dimensional consequence, and whether homeomorphisms must preserve the distinguished cone points are not specified in the supplied record. Completeness and boundary assumptions for the manifolds are also absent. Thus this is a conditional reconstruction of the intended statement, requiring human review; the supplied hypotheses alone are not being claimed here to guarantee existence of tangent cones.
\subsection{Short English statement}
Naber's supplied conjecture says that, outside an exceptional set described as having codimension five, any two tangent cones at a fixed point of a noncollapsed lower-Ricci-curvature limit are homeomorphic. The convergence convention, precise dimension bound, and ambient manifold assumptions require verification. The registry labels the problem open; that status is supplied metadata, not an independently verified assessment.
\subsection{Sources}
\begin{itemize}
\item Local UnsolvedMath registry, permanent ID \(800459\), exact source problem code \texttt{catalogue entry}, title \emph{Spaces with Lower Ricci Curvature --- Conjecture 3}. The supplied record attributes the conjecture to Naber, with proposed year 2020, and labels it open.
\item Supplied source list: Naber, \emph{Conjectures and Open Questions on Spaces with Lower Ricci Curvature} (2020), Conjecture 3. Supplied URL: \url{https://arxiv.org/abs/2010.10031}. The record reports access on 2026-10-16 and extraction by \texttt{arxiv\_latex}. These are supplied provenance details; the source and its mathematical or literature status have not been independently checked here.
\end{itemize}
\end{document}
